\section{Gobernanza y trabajo futuro}

\subsection{Marco de gobernanza}

GR00T no se limita a la investigación, también busca alinearse con la regulación:

\begin{itemize}
    \item \textbf{Safety board}: antes de cada release debe haber una evaluación humana de los reward escalation scenarios
    \item \textbf{Usage policy}: define qué tareas pueden ejecutarse de forma autonomous y cuáles deben ser human-in-loop
    \item \textbf{Audit trail}: registra de forma completa el dataset, los hyperparam y el model ID de cada deployment
\end{itemize}

\begin{importantbox}{Lista de verificación de governance de nivel industrial}
1. ¿Existe un failure mode radical (por ejemplo, tomar un objeto punzocortante)? 2. ¿Se proporciona un override switch? 3. ¿Se registran la instrucción y la response? 4. ¿Existe una alerta de non-nominal behavior?
\end{importantbox}

\subsection{Líneas futuras de investigación}

Jim Fan plantea tres caminos:

\begin{enumerate}
    \item \textbf{LLM + physics solvers}: incorporar differentiable physics a la planeación
    \item \textbf{Self-supervised manipulation}: sintetizar demonstrations de forma automática mediante observation consistency
    \item \textbf{Multi-agent robot swarms}: hacer que múltiples versiones de GR00T colaboren para completar procesos industriales
\end{enumerate}

\subsection{Resumen de la sección}

La gobernanza, la usage policy, el audit trail y las líneas futuras de investigación forman en conjunto la base de seguridad para la implementación industrial.

